% ============================================================
% D65 — Two Limits of the Same Surface: Behavioural Correspondences
%        Between Nucleon and Black Hole Closures in PDL
%
% Author  : Cedric Laubscher
% Corpus  : PDL Document Series, D65
% Depends : D23v2, D29, D33, D37, D42, D43, D44, D50, D54, D63, D64
% Introduces: OP-D65-1, OP-D65-2, OP-D65-3
% Sequel to: D64
% ============================================================

\documentclass[12pt,a4paper]{article}

\usepackage[T1]{fontenc}
\usepackage[utf8]{inputenc}
\usepackage{lmodern}
\usepackage{microtype}
\usepackage{amsmath,amssymb,amsthm}
\usepackage{newpxtext}
\usepackage{mathtools}
\usepackage{booktabs}
\usepackage{array}
\usepackage{longtable}
\usepackage{multirow}
\usepackage{hyperref}
\usepackage[margin=2.5cm]{geometry}
\usepackage{enumitem}
\usepackage{float}
\usepackage[numbers]{natbib}
\usepackage{xcolor}
\usepackage{tikz}
\usetikzlibrary{arrows.meta,positioning,calc}

\definecolor{pdllink}{RGB}{0,60,120}
\definecolor{proofblue}{RGB}{0,70,140}
\definecolor{conjecturegreen}{RGB}{0,120,60}
\definecolor{speculativeorange}{RGB}{180,90,0}
\definecolor{negativered}{RGB}{160,20,20}
\hypersetup{colorlinks=true, linkcolor=pdllink,
citecolor=pdllink, urlcolor=pdllink}

% ---- Theorem environments ------------------------------------------
\newtheorem{theorem}{Theorem}[section]
\newtheorem{lemma}[theorem]{Lemma}
\newtheorem{corollary}[theorem]{Corollary}
\newtheorem{proposition}[theorem]{Proposition}
\newtheorem{definition}[theorem]{Definition}
\newtheorem{remark}[theorem]{Remark}
\newtheorem{openproblem}{Open Problem}
\newtheorem{resolvedproblem}{Resolved Problem}

% ---- coloured boxes, per the N02/D64 protocol -----------------------
\usepackage[most]{tcolorbox}
\newtcolorbox{speculativebox}[1][]{%
colback=orange!5, colframe=speculativeorange, fonttitle=\bfseries, title=Structural correspondence (not a proof of identity), #1}
\newtcolorbox{negativebox}[1][]{%
colback=red!5, colframe=negativered, fonttitle=\bfseries, title=Negative result, #1}
\newtcolorbox{conjecturebox}[1][]{%
colback=green!5, colframe=conjecturegreen, fonttitle=\bfseries, title=Conjecture, #1}

% ---- PDL macros -----------------------------------------------------
\newcommand{\Kfour}{K_4}
\newcommand{\hb}{\hbar}
\newcommand{\PDL}{\textsc{pdl}}
\newcommand{\Rsurf}{R_{\mathrm{surf}}}
\newcommand{\Rtot}{R_{\mathrm{tot}}}
\newcommand{\Rval}{R_{\mathrm{val}}}
\newcommand{\rval}{r_{\mathrm{val}}}
\newcommand{\Rsea}{R_{\mathrm{sea}}}
\newcommand{\Omegasurf}{\Omega_{\mathrm{surf}}}
\newcommand{\Ssurf}{S_{\mathrm{surf}}}
\newcommand{\SBH}{S_{\mathrm{BH}}}
\newcommand{\lPl}{\ell_{\mathrm{Pl}}}
\newcommand{\Meff}{M_{\mathrm{eff}}}
\newcommand{\MPl}{M_{\mathrm{Pl}}}
\newcommand{\sig}{\sigma}
\newcommand{\vp}{\varphi}
\newcommand{\kap}{\kappa}
\newcommand{\Dmiso}{\Delta m_{\mathrm{iso}}}
\newcommand{\GPDL}{G_{\mathrm{PDL}}}
\newcommand{\Geff}{G_{\mathrm{eff}}}
\newcommand{\chif}{\chi}
\newcommand{\Hcycl}{\mathcal{H}_{\mathrm{cycl}}}
\newcommand{\lamC}{\lambda_C}
\newcommand{\Zsat}{Z_{\mathrm{sat}}}

% ============================================================
\begin{document}

\title{\textbf{Two Limits of the Same Surface: Behavioural
Correspondences Between Nucleon and Black Hole Closures in
\PDL}\\[6pt]
{\large \PDL\ Document D65 --- Sequel to D64\\
Three New Open Problems (OP-D65-1, OP-D65-2, OP-D65-3)}}

\author{C\'edric Laubscher\\[4pt]
\small Independent Researcher, Switzerland\\
\small \href{https://cedriclaubscher.ch}{cedriclaubscher.ch}}

\date{July 2026}

\maketitle

% ============================================================
\begin{abstract}
Document D64 established a structural correspondence between the
proton's surface degeneracy theorem (D50) and the soft-hair
programme of Hawking, Perry and Strominger, and reported that a
naive nucleon-counting extrapolation of the \PDL\ per-relation
degeneracy fails the Bekenstein--Hawking value by 36.7 orders of
magnitude, independently of black hole mass. This document is a
direct sequel. It does not attempt to close that gap. Instead it
asks a different, prior question: independently of any absolute
microstate count, does \PDL\ predict shared \emph{behaviours}
between its two most extreme closures --- the isolated proton and
the macroscopic black hole horizon --- that are not visible in
either limit alone? Five such behaviours are established or
substantially developed here. First, both closures are shown to
require, as an unconditional consequence of $\Kfour$'s
$SL(2,\mathbb{C})$ pulsation structure (D33) rather than of any
continuum approximation, an internal state space $\Hcycl$ richer
than a binary sign; applied to a closed spherical \PDL\ surface,
this forces exactly two coherence defects, never zero, by a
combinatorial analogue of the Poincar\'e--Hopf theorem verified
independently by spinor relaxation, exact discrete Gauss--Bonnet
computation, and direct holonomy calculation in arbitrary
precision arithmetic. This gives a structural, axiom-level account
of why no black hole has ever been observed with exactly zero
spin, complementing rather than duplicating the Kerr/\textsc{cft}
correspondence of Guica, Hartman, Song and Strominger, whose
near-horizon $AdS_2\times S^2$ geometry shares the same $S^2$ that
already governs the topological origin of the exponent 18 at the
proton scale (D23 v2). Second, the curvature-migration law
governing the approach to extremality is computed exactly and
shown to have asymptotic exponent 1, not 2, only once floating-point
noise is eliminated by arbitrary-precision arithmetic; this
establishes a genuine tension with the third law of black hole
thermodynamics that is reported openly rather than resolved.
Third, a dissolution threshold is derived from the previously
established nucleon saturation number $\Zsat=20$ (D22) combined
with the standard nuclear force range, giving a critical mass
$M^*\approx4.3\,M_\odot$ at which the density required for
nucleon dissolution coincides exactly with the Schwarzschild
condition; this value falls inside the observed compact-object
mass gap without having been tuned to do so. Fourth, seven
independent attempts to derive an absolute microstate count or a
Bardeen--Thorne-type spin-limiting mechanism from \PDL\ combinatorics
are reported as failures, each with its precise structural cause,
extending and reinforcing the negative result of D64 rather than
contradicting it. Fifth, these results are organised under a single
methodological principle: at the black hole scale, \PDL\ currently
supports claims about \emph{behaviour} (what is topologically
forced, what functional form a law must take, what threshold a
transition must respect) but not about \emph{magnitude} (absolute
entropy, absolute coupling strength), and this asymmetry is shown to
be structural rather than a temporary gap in derivation effort.
Three open problems are introduced.
\end{abstract}

\clearpage
\tableofcontents
\newpage

% ============================================================
\section{Purpose and intended audience}
\label{sec:purpose}
% ============================================================

This document, like D64, is written for colleagues outside the
\PDL\ programme who are familiar with black hole thermodynamics,
the Kerr/\textsc{cft} and soft-hair literature, and standard
compact-object astrophysics, but not with the internal
combinatorics of \PDL. Section~\ref{sec:prereq} recalls, without
proof, exactly the prerequisite \PDL\ results; full proofs are in
the cited documents. Section~\ref{sec:consolidation} briefly
consolidates the negative result of D64 using independent
methods. Section~\ref{sec:topology} presents the principal new
result: the topological necessity of a non-zero spin defect on any
closed \PDL\ surface. Section~\ref{sec:dynamics} reports the
curvature-migration dynamics near extremality, including an
unresolved tension with the third law. Section~\ref{sec:threshold}
derives the dissolution threshold and critical mass.
Section~\ref{sec:negative} documents, precisely, every mechanism
tested and excluded. Section~\ref{sec:principle} states the
behaviour-versus-magnitude principle. Section~\ref{sec:open}
introduces three new open problems. As with D64: no result in this
document modifies D37, D50, D43 or D44; it is a correspondence and
consolidation document, not a re-derivation of established
theorems.

% ============================================================
\section{PDL prerequisites}
\label{sec:prereq}
% ============================================================

The \PDL\ programme derives physical structure from four
combinatorial axioms (C1--C4) on finite signed graphs, with a
single irreducible external parameter, $\Dmiso=m_d-m_u$, at the
\PDL--\textsc{qcd} interface~\citep{PDL_D01_2026}. C1 requires an
irreducible binary pulsation ($s\to-s$); C2 requires triangular
coherence; C3 requires minimal self-referential closure; C4
requires maximisation of a coherence-weighted functional $\Omega$.
The electron prototype is $\Kfour$, the complete graph on four
vertices and six edges, the smallest closure satisfying C1--C3. The
proton is the unique hierarchical composite satisfying C1--C4,
characterised by the integer quintuplet
$(n_u,n_d,\rval,\Rsea,\Rtot)=(24,28,930,10087,11017)$.

\begin{definition}[Active surface and leakage fraction]
$\Rsurf=310\vp\in\mathbb{Q}(\sqrt5)$ is the subset of the proton's
total relational budget forming stable $(A)\wedge(B)$ couplings
with an external $\Kfour$ closure~\citep{PDL_D29_2026}. The leakage
fraction $\kap=\Rsurf/\Rtot\approx4.6\%$ is the fraction of the
proton's structure that must remain externally coupled rather than
frozen.
\end{definition}

\begin{theorem}[Surface locality and surface degeneracy,
D37 and D50]
\label{thm:BH1D50}
All entropy of an isolated \PDL\ closure is localised on its active
surface: $\Omega_{\mathrm{val}}=1$, $\Omegasurf=4^{\Rsurf}$, an
unconditional theorem obtained by exhaustive enumeration of the
$8\times6\times16=768$ configurations of the $(A)\wedge(B)$
stability criterion~\citep{PDL_D29_2026,PDL_D37_2026,PDL_D50_2026}.
\end{theorem}

\begin{theorem}[Dissolution of hierarchical structure, D63]
\label{thm:D63}
Valence cores are permanent partial closures with an open boundary
requiring coupling to a relational sea; they are not
self-contained. No autonomous $\Kfour$-type closure has been
formalised for a deconfined valence core once the hosting sea is
destroyed~\citep{PDL_D63_2026}.
\end{theorem}

\begin{theorem}[Topological origin of the exponent 18, D23 v2]
\label{thm:D23}
$18=6+5+4+3$ is derived as a topological necessity from
$H_2(S^2;\mathbb{Z})\cong\mathbb{Z}$, imposing a rank deficit of
exactly one dimension at each transition of the chain
$C_0\to C_1\to C_2\to C_3$ (proton $\to$ nucleus $\to$
matter/gravitation), recovered exactly by three boundary
morphisms~\citep{PDL_D23v2_2026}. This is the origin of the
exponent in $\GPDL=(\hb c/m_p^2)\varepsilon_G^{18}$
\citep{PDL_D43_2026,PDL_D44_2026}.
\end{theorem}

\begin{definition}[Pulsation structure of $\Kfour$, D33]
\label{def:D33}
Every $\Kfour$ closure carries a two-dimensional complex state
space $\Hcycl$ and a unique half-cycle operator $T=-i\tau_2\in
SU(2)$ satisfying $T^2=-\mathbb{1}$, $T^4=\mathbb{1}$, at Compton
frequency $\omega_c=m_pc^2/\hb$~\citep{PDL_D33_2026}. Uniqueness is
a theorem: no other operator is compatible with C1--C4.
\end{definition}

\begin{theorem}[Nucleon saturation, D22]
\label{thm:Zsat}
A neutron's simultaneous binding capacity to protons saturates at
$\Zsat=\lfloor T/(T-T_{pp})\rfloor+1=20$; beyond this, additional
engagement adds no further binding, $C(Z>\Zsat)=190\,T_{pp}=$
constant~\citep{D22_nuclear_stability}.
\end{theorem}

\begin{theorem}[Density-dependent gravitational coupling, Gate 3]
\label{thm:Gate3}
$\Geff(N)=\sig(N)\GPDL$, $\sig(N)=1-(1-\kap)^N$, proven exactly in
the dilute regime $N\kap\ll1$; the extension to $N\kap\gg1$ (the
regime relevant to macroscopic closures) is not established.
\end{theorem}

% ============================================================
\section{Consolidating the negative result of D64}
\label{sec:consolidation}
% ============================================================

D64 reported that a naive nucleon-pairing extrapolation of
$\Omegasurf$ fails the Bekenstein--Hawking entropy by 36.7 orders
of magnitude, independently of black hole mass, verified across
four masses spanning eight orders of magnitude. Four further
attempts, developed independently of D64 and using different
physical reasoning in each case, converge on discrepancies in the
same range.

\begin{negativebox}
Four independent constructions for a macroscopic entropy count,
none aware of the specific numerical target in advance:
\begin{enumerate}[nosep]
\item Sequential accretion chain (bounded active-surface width
$W\approx3$): discrepancy $\sim10^{-20}$ (undershoot).
\item Each macroscopic site treated as one isolated proton's full
$\Rsurf$: discrepancy $\sim10^{+40}$ (overshoot).
\item Saturated boundary shell at $\Zsat=20$: discrepancy
$\sim10^{+38.6}$ (overshoot).
\item Elementary area calibrated once at the proton's Compton
wavelength $\lamC$, applied uniformly to the horizon area:
discrepancy $\sim10^{-35.9}$ (undershoot).
\end{enumerate}
All four discrepancies are of the same order as
$\alpha_G=\varepsilon_G^{18}\sim10^{-39}$, the dimensionless
gravitational coupling already established at the proton scale
(D43, D44). This is shown, not merely observed: for a fixed
Planck-length-based elementary area, the ratio of any
particle-physics length scale squared to $\lPl^2$ is algebraically
forced to be of order $\alpha_G^{-1}$, since $\lPl^2=\hb G/c^3$
already contains $G$. No test of this form can be a genuine
independent derivation of $G$ or of $\SBH$; each is a restatement
of the hierarchy problem, not a \PDL-specific failure.
\end{negativebox}

\begin{remark}
A parallel diagnosis applies to the combinatorial size variable
itself. Any astronomical mass $M$ in kilograms is obtained from the
directly measured, model-independent product $GM$ by dividing by a
separately measured value of $G$; a mass so obtained already
contains $G$. $N=M/m_p$ is therefore circular for the same reason
$\lPl$ is circular as a probe of $G$: both import the very quantity
under test through the back door of unit conversion.
\end{remark}

% ============================================================
\section{Topological forcing of a non-zero spin defect}
\label{sec:topology}
% ============================================================

This section departs from D64 in subject matter entirely: D64
concerns entropy counting; this section concerns spin, and is, to
the author's knowledge, the first attempt within \PDL\ to derive a
structural constraint on black hole angular momentum from C1--C4
rather than from imported general-relativistic geometry.

\subsection{Why a binary sign is insufficient}

\begin{negativebox}
A first attempt modelled the coherence field on a closed \PDL\
surface as a binary sign $\sig_v\in\{+1,-1\}$ per vertex (C1 taken
alone), with edge signs $s_{ij}=\sig_i\sig_j$. For any such
assignment, on any graph whatsoever,
$s_{ij}s_{jk}s_{ik}=\sig_i^2\sig_j^2\sig_k^2=1$ identically: C2
(triangular coherence) is always trivially satisfiable. This is
confirmed by general signed-graph theory (the frustration index of
a $\mathbb{Z}_2$-signed graph is a property of the graph alone,
never forced by the topology of an embedding surface). A richer
state space is therefore required to see any topological
obstruction at all.
\end{negativebox}

\subsection{\texorpdfstring{The $SU(2)$ obstruction}{The SU(2) obstruction}}

Definition~\ref{def:D33} supplies exactly the required richness: an
element of $\Hcycl\cong\mathbb{C}^2$ per vertex, with nearest-neighbour
transport given by the spinorial lift of the discrete
transport angle $\theta_{ij}$ (the standard holonomy representation
of a rotation in the local tangent plane, $U(\theta)=\exp(i\theta
\sigma_z/2)$).

\begin{theorem}[Forced coherence defects on a closed surface]
\label{thm:2defects}
For any \PDL\ surface of spherical topology, discretised as a
closed triangulation with the $\Hcycl$ transport structure above, a
minimum of exactly two points at which the local coherence
condition cannot be satisfied is unavoidable: no configuration with
zero or one such point exists.
\end{theorem}

This was verified by three independent methods, in order of
increasing rigour.

\paragraph{Method 1 --- spinor relaxation.} A Ginzburg-Landau-type
relaxation of the $\Hcycl$-valued field, with the $(A)\wedge(B)$-
motivated transport structure, was run on an icosphere ($642$
vertices) and on Fibonacci-sphere meshes ($400$--$2000$ vertices),
each with 20 independent random initial seeds. The minimum number
of amplitude zeros (coherence defects) obtained was $2$ on every
run and every mesh; it was never $0$ or $1$; occasional excess ($4$,
$6$) corresponds to kinetically trapped, unannihilated
vortex--antivortex pairs, confirmed by re-running with longer
relaxation. A flat, topologically trivial control (a periodic
square lattice with zero curvature) reached exactly $0$ defects on
every run, confirming the obstruction is topological, not an
artefact of the relaxation method.

\paragraph{Method 2 --- exact discrete Gauss--Bonnet.} Independently
of any relaxation dynamics, the discrete angular deficit
$\delta_v=2\pi-\sum(\text{face angles at }v)$ was computed exactly.
For every closed triangulation tested (icosahedron, $12$ vertices;
Fibonacci spheres, $400$ and $2000$ vertices), $\sum_v\delta_v=4\pi$
to machine precision, a pure consequence of $V-E+F=2$, requiring no
relaxation and no approximation. An open triangulated patch, by
contrast, carries zero deficit at every interior vertex, with the
full deficit absorbed at the boundary alone.

\paragraph{Method 3 --- direct holonomy.} The spinor holonomy
$U=\exp(i\delta_v\sig_z/2)$ was evaluated directly at every vertex
from the exact deficit of Method 2. Every vertex with $\delta_v\neq0$
was confirmed to carry non-trivial holonomy ($U\neq\mathbb{1}$),
matching the defect count of Method 2 exactly, vertex for vertex,
on every mesh tested.

\begin{remark}
The three methods are logically independent (a physically motivated
relaxation, a pure combinatorial identity, and an exact matrix
calculation) and converge on the identical count. This is offered
as the paper's most robust single result.
\end{remark}

\subsection{Where the defects sit, and why this is spin}

At fixed oblateness parameter $\chif$ (the flattening
$s=\sqrt{1-\chif^2}$, chosen to match the separation of the two
Kerr horizon radii, $s=(r_+-r_-)/(2GM/c^2)$), 25 independent
relaxation runs place the two coherence defects at the poles of the
imposed flattening axis, essentially uniquely (a single minimum,
not a degenerate family), in every run. The axis of flattening is
not an external input to be justified separately: it is, on this
evidence, the locus that the coherence field itself is forced to
select as the site of its two unavoidable defects. This gives a
structural, axiom-level candidate explanation for the empirical
fact that no black hole has ever been measured with $\chif=0$.

\begin{speculativebox}
Guica, Hartman, Song and Strominger showed that the near-horizon
geometry of an extremal Kerr black hole factorises as
$AdS_2\times S^2$~\citep{GuicaHartmanSongStrominger2009}, with the
entropy recovered from a chiral Virasoro algebra associated with
the axial isometry. The $S^2$ appearing there is not asserted here
to be mathematically identical to the $S^2$ of
Theorem~\ref{thm:D23} (the topological origin of the exponent 18)
or of Theorem~\ref{thm:2defects}; the three results share no direct
derivation. What is shared is the topological type: the same
two-sphere, entering three independently-motivated arguments
(a hierarchical filtration exponent at the proton scale, a
combinatorial spin obstruction at the horizon scale, and a
near-horizon isometry in general relativity), in each case as the
source of a non-trivial physical consequence rather than as inert
background. Whether this is more than a shared topological
vocabulary is an open question, stated formally as OP-D65-1 below.
\end{speculativebox}

% ============================================================
\section{Dynamics near extremality: an open tension}
\label{sec:dynamics}
% ============================================================

\subsection{The curvature-migration law}

Tracking the exact (arbitrary-precision, 50-digit) angular deficit
of a fixed set of near-pole vertices as $\chif$ is increased shows
the total deficit at the poles falling monotonically, migrating
towards the equatorial band, while the total ($4\pi$) is conserved
throughout.

\begin{theorem}[Asymptotic exponent]
\label{thm:exponent}
The polar deficit obeys $\delta_{\mathrm{pole}}(\chif)\to
a(1-\chif^2)$ as $\chif\to1$, with local log--log exponent
converging to $1.000000$ to six decimal places for
$\chif\geq0.995$, computed in 50-digit arbitrary precision. Standard
double-precision (float64) computation of the same quantity gives a
spurious exponent near $2$, an artefact of catastrophic cancellation
at deficit values below $10^{-6}$; the discrepancy is resolved
entirely by precision, not by any change of model.
\end{theorem}

\begin{negativebox}
This directly contradicts an earlier working hypothesis (developed
in the same investigation and reported here for completeness) that
the migration rate should vanish smoothly at $\chif=1$, mirroring
the vanishing surface gravity of extremal Kerr. Fitting the
migration rate $f(\chif)=-d\delta_{\mathrm{pole}}/d\chif$ against
three candidate forms --- a free power law, a form constrained to
diverge as $(1-\chif^2)^{-1/2}$ (motivated by the vanishing surface
gravity), and an unconstrained pure power law --- shows the
constrained, physically motivated form to be the \emph{worst} fit
by a corrected Akaike information criterion; the simplest,
unmotivated power law is statistically preferred. $f(\chif)$
saturates to a finite, non-zero value as $\chif\to1$, not zero.
\end{negativebox}

\subsection{Tension with the third law}

Converting the resulting combinatorial evolution parameter $\tau$ to
physical time using the anchor $\tau_{\mathrm{horizon}}=GM/c^3$
(the horizon's own light-crossing time, an invariant computable
directly from the measured product $GM$, in contrast to the
proton's own Compton period $\tau_0=\hb/(m_pc^2)$, which gives a
physically absurd result of order $10^{-25}\,$s for the same
process) yields $\tau_{\mathrm{physical}}\approx0.2\,$ms for a
GW250114-scale merger, the correct order of magnitude for a real
ringdown. However, because $f(\chif)$ remains bounded away from zero
as $\chif\to1$ (Theorem~\ref{thm:exponent}), the combinatorial
integral $\int d\chif/f(\chif)$ converges: this specific
mechanism predicts that exact extremality is reached in \emph{finite}
time, in direct tension with the third law of black hole
thermodynamics~\citep{BardeenCarterHawking1973}, which forbids this
in any finite number of physical processes.

\begin{remark}
This tension is reported, not resolved. Two resolutions are logically
available and are distinguished only by further work: either this
specific topological mechanism genuinely does not reproduce the
third law (a distinctive, falsifiable departure), or the mapping
from combinatorial $\tau$ to physical time is itself divergent near
extremality in a way not captured by the single anchor
$\tau_{\mathrm{horizon}}$ used here. No claim is made in favour of
either option.
\end{remark}

\subsection{\texorpdfstring{A partially motivated numerical value for $\chif$}{A partially motivated numerical value for chi}}

\begin{conjecturebox}
Solving the self-consistency condition $\chif=\beta(\chif)/2$, with
$\beta(\chif)=1+D(\chif)$ and $D(\chif)=1-\sqrt{1-\chif^2}$ the
normalised separation of the two Kerr horizon radii (interpreted as
an inertia-like shape factor going from $1$, an equatorial ring, at
$\chif=0$, to $2$, exact degeneracy of the two roots, at
$\chif=1$), gives $\chif=3/5$ exactly. This value differs from the
measured GW250114 remnant spin ($\chif=0.68$) by $12\%$, and from
GW150914 ($\chif=0.67$) by a comparable amount. Two independently
motivated forms of $\beta(\chif)$ (the one above, and one built from
the ratio $r_-/r_+$ instead of $D$) both admit non-trivial
self-consistent solutions, clustering at $\chif=0.54$--$0.60$;
unmotivated polynomial forms with the same boundary values admit no
non-trivial solution at all. A candidate mechanistic justification
(the inner horizon contributing to angular momentum with a weight
equal to the outer horizon's, once the two roots coincide) was
tested against the established two-horizon thermodynamic literature
and found inconsistent: angular momentum is the same conserved
quantity in the first law of both the outer and inner horizon,
never partitioned between them. The value $\chif=3/5$ is therefore
reported as a reproducible numerical pattern with a partial,
geometrically motivated derivation, explicitly \emph{not} as a
theorem.
\end{conjecturebox}

An out-of-sample test against Cygnus~X-1 ($\chif>0.998$,
\citep{Zhao2021}) fails by a wide margin, as expected: this object
is spun up by sustained accretion over millions of years, a
mechanism entirely absent from the present topological
construction, which concerns only the state at closure formation.
The failure is diagnostic rather than damaging: it delimits the
domain of the mechanism (freshly formed closures) rather than
falsifying it.

% ============================================================
\section{The dissolution threshold and a critical mass}
\label{sec:threshold}
% ============================================================

\subsection{A density threshold, independent of total size}

\begin{theorem}[Local dissolution threshold]
\label{thm:threshold}
Modelling local nucleon coupling using $\Zsat=20$
(Theorem~\ref{thm:Zsat}) as the maximum number of simultaneously
engageable neighbours within the standard nuclear coupling range
$r_c$, the mean local neighbour count in the interior of a
uniform-density cluster depends only on the local number density
$\rho$, never on the total cluster size $N$, once boundary effects
are correctly excluded and Poisson statistics are properly accounted
for (confirmed for $N=200$ to $12\,800$, with the observed
dissolved fraction matching the Poisson-theoretic prediction to
within sampling noise at every density tested). The threshold
density is
\begin{equation}
\rho^*=\frac{\Zsat}{\frac{4}{3}\pi r_c^3}.
\end{equation}
For $r_c=2\,$fm, $\rho^*=3.73\,\rho_0$, where $\rho_0$ is ordinary
nuclear saturation density.
\end{theorem}

\subsection{A critical mass}

Because the radius at fixed density scales as $M^{1/3}$ while the
Schwarzschild radius scales as $M$, there exists exactly one mass
$M^*$ at which these coincide:
\begin{equation}
M^*=\frac{c^3}{\sqrt{\tfrac{32}{3}\pi\rho^*G^3}}\approx4.31\,M_\odot.
\end{equation}

\begin{remark}
For $M<M^*$, the dissolution radius exceeds the Schwarzschild
radius: a collapsing object can, in principle, reach $\rho^*$ while
still larger than its own horizon, permitting a stable dissolved
(single coherent surface) configuration that is not a black hole,
supported against further collapse by physics not yet derived here
(Section~\ref{sec:negative}). For $M>M^*$, the Schwarzschild radius
is reached first, at ordinary density; dissolution, if it occurs at
all, happens later, already hidden behind the horizon.
\end{remark}

\begin{speculativebox}
$M^*\approx4.3\,M_\odot$ falls inside the observed compact-object
mass gap, the interval between the heaviest well-measured neutron
stars ($\sim2.2$--$2.5\,M_\odot$) and the lightest black holes
traditionally identified by dynamical mass measurement
($\sim5\,M_\odot$), an interval containing at least one
observationally ambiguous object, the $2.6\,M_\odot$ component of
GW190814~\citep{GW190814}. A sensitivity scan over
$r_c\in[1.0,3.0]\,$fm shows $M^*$ spans the interval
$[2.5,5.0]\,M_\odot$ for $r_c\in[1.39,2.21]\,$fm, a range
overlapping the standard pion Compton wavelength
($\hbar/m_\pi c\approx1.4\,$fm), often cited as the natural
one-pion-exchange range. Neither the mass-gap interval nor the
pion-range coincidence was a target of the construction; both
values are outputs, not inputs.
\end{speculativebox}

% ============================================================
\section{What was tried and excluded}
\label{sec:negative}
% ============================================================

The following mechanisms were formulated as candidate explanations
for the mechanism supporting a sub-$M^*$ dissolved object against
further collapse, or for a Bardeen--Thorne-type spin-limiting
mechanism at $\chif\approx0.998$, and are reported here as excluded,
with the precise reason in each case, following the same discipline
applied throughout the \PDL\ corpus to negative results.

\begin{negativebox}
\textbf{Rigid-body inertia for Kerr.} Requiring $I=\beta MR^2$ with
$\Omega_H=c^3/(2GM)$ (the extremal horizon angular velocity) to
reproduce $J_{\mathrm{ext}}=GM^2/c$ exactly forces $\beta=2$,
exceeding the absolute maximum ($\beta=1$, a thin equatorial ring)
for any mass distribution confined within radius $R$. The classical
rigid-rotation picture is not merely miscalibrated; it is
kinematically impossible for this quantity.

\textbf{Two-horizon thermodynamics as the origin of the spin
doubling.} The candidate picture in which the inner and outer
horizons each contribute independently to angular momentum,
summing at exact degeneracy, is inconsistent with the established
first law of inner-horizon thermodynamics: $J$ is the identical
conserved quantity in both the outer and inner first laws, never
split between them.

\textbf{Geometric saturation of local curvature.} The maximum local
angular deficit at the most-loaded equatorial vertex, computed in
arbitrary precision, remains below $11\%$ of its absolute geometric
bound ($\pi$, for a degree-3 vertex) even at $\chif=0.9999$: no
local saturation mechanism intervenes before exact extremality on
any mesh tested.

\textbf{Spin stress-energy suppression (D48v3).} The known
suppression $C_{\mathrm{spin}}\sim(\lamC/L)^2\,C_{\mathrm{mass}}$,
evaluated at $L=r_+$ for a real black hole, is non-zero (the
homogeneous-limit theorem $C_{\mathrm{spin}}=0$ applies only as
$L\to\infty$) but of order $10^{-42}$, and varies by at most a
bounded factor of 4 between $\chif=0$ and $\chif=1$ (since $r_+$ can
only shrink by a factor of 2): far too small and far too gently
varying to act as a brake.

\textbf{Coincidence of the 126 nuclear magic number.} The
near-identity $126\times(\Delta n+1)^2\approx2\pi\Rsurf(p)$
(agreement $0.05\%$), initially considered a candidate holonomy
threshold, is a comparison between a physical nucleon count (126,
an entity count) and a \PDL\ relational quantity; on the same
grounds established throughout this document for rejecting
$N=M/m_p$, this comparison mixes categories and is not retained as
a structural result.

\textbf{Polytropic pressure support, $\Gamma=5/3$.} The scaling
$P\sim(\lamC/\ell)^2\rho c^2$ with local interparticle spacing
$\ell(\rho)=(m_n/\rho)^{1/3}$ produces exactly a $\Gamma=5/3$
polytropic equation of state; explicit Newtonian stellar-structure
integration shows this equation of state admits no maximum mass at
all ($M\propto\sqrt{\rho_c}$, unbounded), in contrast to the
$\Gamma=4/3$ relativistic degenerate equation of state required for
a genuine Chandrasekhar-type limit. An apparent numerical
coincidence with $M^*$ at one particular central density was
verified, on extending the density scan, to be an artefact of where
the scan was stopped, not a genuine limit.

\textbf{Correlation between spin and horizon scale.} A four-event
correlation between $\chif$ and $\log_{10}(GM_f m_p/c\hb)$
($r=0.94$) was found not to survive extension to seven
well-measured events of ordinary formation channel ($r=-0.47$,
sign-reversed), an explicit illustration of small-sample artefact
rather than of a genuine trend.
\end{negativebox}

\begin{remark}
Three of these seven failures (geometric saturation, spin
suppression, and the equation of state) share a common structural
cause, stated formally in Section~\ref{sec:open} as OP-D65-3: every
currently rigorous \PDL\ result invoked here (the pressureless
equation of state of D54, the vanishing spin tensor of D48v3, the
absolute geometric bound on local curvature) is proven only in the
homogeneous or isolated-closure limit. None has yet been derived
for a macroscopically structured, symmetry-broken closure --- the
regime a black hole horizon actually occupies.
\end{remark}

% ============================================================
\section{Behaviour versus magnitude}
\label{sec:principle}
% ============================================================

The results and failures above organise cleanly under a single
principle, stated here as the paper's central methodological claim
rather than as a further technical result.

\begin{conjecturebox}[title=Behaviour-versus-magnitude principle]
At the black hole scale, every \PDL\ result that has survived
independent verification concerns a \emph{behaviour}: what
configuration is topologically permitted or forbidden
(Theorem~\ref{thm:2defects}), what functional form a law must take
near a boundary (Theorem~\ref{thm:exponent}), what threshold a
phase transition must respect (Theorem~\ref{thm:threshold}), or how
a known quantity scales with a control parameter
($R_s$ and $T_H$ depend only on $GM$, never on $G$ and $M$
separately). Every attempt to instead predict an absolute
\emph{magnitude} at the black hole scale --- an entropy value, a
coupling strength, a maximum supportable mass from first principles
--- has failed, and failed for a reason traceable, in every case
examined, either to the hidden circularity of importing $G$ through
a length or mass conversion, or to the unresolved extension of
homogeneous-regime theorems to the inhomogeneous, symmetry-broken
regime a real horizon occupies. This is proposed as a structural
property of the current theory, delimiting its domain of validity
at the black hole scale, not as a temporary shortage of successful
calculations.
\end{conjecturebox}

% ============================================================
\section{Open problems}
\label{sec:open}
% ============================================================

\begin{openproblem}[The shared $S^2$]
\label{OP-D65-1}
Is the topological type $S^2$, appearing independently in the
origin of the exponent 18 (Theorem~\ref{thm:D23}, proton scale), in
the forced spin defect (Theorem~\ref{thm:2defects}, horizon scale),
and in the near-horizon $AdS_2\times S^2$ geometry of extremal
Kerr~\citep{GuicaHartmanSongStrominger2009}, more than a shared
topological vocabulary --- specifically, does a single combinatorial
construction on $\Kfour$ exist from which all three follow as
corollaries, rather than as three independently-derived
consequences of the same manifold type? Entry point: D23~v2,
Section~\ref{sec:topology} of this document,
\citep{GuicaHartmanSongStrominger2009,HackoHawkingPerryStrominger2018}.
\end{openproblem}

\begin{openproblem}[Combinatorial-to-physical time near extremality]
\label{OP-D65-2}
Does the tension identified in Section~\ref{sec:dynamics}
(finite combinatorial time to reach $\chif=1$, in tension with the
third law) resolve because the mapping from combinatorial $\tau$
to physical time itself diverges near extremality, or does it
reflect a genuine, falsifiable departure of this specific \PDL\
mechanism from the third law? A resolution requires either an
independent, non-circular derivation of the $\tau\to$ physical-time
map beyond the single anchor $GM/c^3$ used here, or an explicit
demonstration that no such divergent map exists, which would
constitute a distinctive prediction. Entry point:
Section~\ref{sec:dynamics}, \citep{BardeenCarterHawking1973}.
\end{openproblem}

\begin{openproblem}[A \PDL\ theory of the inhomogeneous, symmetry-broken
regime]
\label{OP-D65-3}
Construct the coherence stress-energy tensor, the equation of
state, and the local curvature bound of D48v3 and D54 for a
macroscopically structured closure with a preferred axis (finite
inter-cell coherence-variation scale $L$), rather than in the
homogeneous limit $L\to\infty$ on which every current theorem in
that framework depends. This is identified, in
Section~\ref{sec:negative}, as the common cause of three
independent negative results in this document, and is proposed as
the single most consequential open problem for future \PDL\ work on
black holes, superseding the separate pursuit of any one of its
symptoms. Entry point: D48~v3, D54, Remark following
Section~\ref{sec:negative}.
\end{openproblem}

% ============================================================
\section{Conclusion}
\label{sec:conclusion}
% ============================================================

D64 asked whether a single microstate-counting theorem, established
for an isolated proton, could be extrapolated to a macroscopic
black hole, and found that it could not, by a precisely
characterised 36.7 orders of magnitude. This document asked a
narrower but, we think, more tractable question: independently of
that gap, does the same set of four axioms impose shared
\emph{behaviours} on both limits? The answer developed here is a
qualified yes. A structural, topologically forced account of why
black hole spin is never exactly zero has been established with a
rigour --- three independent, convergent, arbitrary-precision
methods --- comparable to the best-verified results elsewhere in
the \PDL\ corpus. A density threshold derived from a
nuclear-physics theorem established years before this
investigation began reproduces, without adjustment, the order of
magnitude and approximate location of an observed astrophysical
mass gap. Seven candidate mechanisms for extending these results
towards absolute magnitudes were tested and honestly excluded, each
for a distinct, stated reason, converging on a single deeper open
problem: the absence, throughout \PDL, of any theorem for the
inhomogeneous regime a real horizon occupies. We consider this
convergence, rather than any single positive result, to be the
most useful outcome of the investigation reported here.

% ============================================================
\bibliographystyle{unsrt}
\bibliography{D65_references}

\end{document}